\section{Convex fitting and a numerical optimality certificate}
\label{app:optimization}

This appendix describes the fitted decoder, including a computable
regularized-objective gap.  The gap concerns the finite training objective;
it is distinct from the population prediction-risk certificate.

Let $F\in\mathbb R^{N\times P}$ contain the frozen geometric feature
rows and $Y\in\mathbb R^{N\times H}$ the corresponding scalar targets
at $H$ horizons.  Different output coordinates can also be placed in
different columns.  The training windows can be dependent: optimization
of this finite objective makes no independence claim about its rows.
With an unpenalized intercept $a\in\mathbb R^H$, the fitted coefficient
matrix solves
\begin{equation}
 \min_{W,a}\frac1N\|Y-\boldsymbol1a^\top-FW\|_{\rm F}^2
                       +\lambda\|W\|_{\rm F}^2
 \quad\text{subject to}\quad
       \sum_{b=1}^s\|W_b\|_{\rm F}\leq\kappa.
 \label{eq:training-program}
\end{equation}
Here $W_b$ contains the rows belonging to block $b$ and $\lambda\geq0$.
The loss sums the horizon losses before averaging over rows; it is not
additionally divided by $H$.  Clipping is applied after fitting.  Thus
\eqref{eq:training-program} is a convex surrogate for the evaluated
clipped-forecast risk, not an assertion that the clipped empirical loss
was globally minimized.

Eliminating the intercept gives $a^\top=\bar Y-\bar F W$.  With centered
matrices $F_c,Y_c$, write
\[
 M=F_c^\top F_c/N+\lambda I,
 \qquad C=F_c^\top Y_c/N,
 \qquad c=\|Y_c\|_{\rm F}^2/N.
\]
The reduced objective and gradient are
\[
 J(W)=\operatorname{tr}(W^\top M W)
             -2\langle W,C\rangle_{\rm F}+c,
 \qquad G(W)=2(MW-C).
\]
Its gradient is Lipschitz in Frobenius norm with constant $2\|M\|_2$.
The implementation uses projected gradient steps with acceleration and
monotone restarts.  The following formulas allow feasibility and
optimization quality to be assessed independently of convergence-rate
claims about that particular implementation.

\begin{proposition}[Projection onto the block norm ball]
\label{prop:block-projection}
Let $A$ have the same block partition as $W$ and define
$a_b=\|A_b\|_{\rm F}$.  If $\sum_ba_b\leq\kappa$, its Euclidean
projection onto the feasible set in \eqref{eq:training-program} is $A$.
Otherwise, let $\tau\geq0$ satisfy
$\sum_b(a_b-\tau)_+=\kappa$.  The projection is
\begin{equation}
 \bigl[\mathcal P_\kappa(A)\bigr]_b
 =\begin{cases}
       (1-\tau/a_b)_+ A_b,& a_b>0,\\
       0,& a_b=0.
   \end{cases}
 \label{eq:block-projection}
\end{equation}
When $\kappa=0$ the projection is the zero matrix.  For a feasible $W$, a
regularized-objective suboptimality bound is
\begin{equation}
 0\leq J(W)-J^*
 \leq \mathcal D(W)
 :=\langle G(W),W\rangle_{\rm F}
                    +\kappa\max_b\|G_b(W)\|_{\rm F},
 \label{eq:optimization-gap}
\end{equation}
where $J^*$ is the minimum of \eqref{eq:training-program}.
\end{proposition}

\begin{proof}
For a fixed block norm $v_b\geq0$, Cauchy--Schwarz shows that the
closest block to $A_b$ points in its direction and has squared distance
$(v_b-a_b)^2$.  The projection therefore reduces to Euclidean projection
of $(a_b)_b$ onto $\{v_b\geq0:\sum_bv_b\leq\kappa\}$.
The scalar Karush--Kuhn--Tucker conditions give
$v_b=(a_b-\tau)_+$ and the stated threshold equation, proving
\eqref{eq:block-projection}.  For the gap, convexity gives
$J(W)-J(V)\leq\langle G(W),W-V\rangle_{
m F}$ for every feasible $V$.
The minimum of $\langle G(W),V\rangle_{
m F}$ over the block norm
ball is $-\kappa\max_b\|G_b(W)\|_{
m F}$: put the entire budget in a
block of largest gradient norm and point opposite to that gradient.
Substituting an optimal $V$ proves \eqref{eq:optimization-gap}.
\end{proof}

The implementation records the coefficient norms, objective gap,
iteration count, and whether the requested gap tolerance was reached.
These are numerical quantities evaluated in floating point; they are
not interval-arithmetic certificates.  Prediction-risk validity only
requires that the returned model satisfy the Lipschitz budget used in
calibration.  It does not require a globally optimal training solution.
For the joint multi-horizon program, each column satisfies
$\sum_b\|(W_b)_{\cdot h}\|_2\leq
\sum_b\|W_b\|_{
m F}\leq\kappa$, so the same declared budget is
valid at every horizon.  The objective's squared Frobenius ridge term
is not invariant to duplicated feature blocks, even though the
constraint's predictor class is invariant.  The unregularized
duplicate-block test therefore sets $\lambda=0$.

